\documentclass[lang=cn,color=blue,math=cm,device=normal]{elegantbook}

\usepackage{array}
\usepackage{longtable}
\usepackage{tabularx}
\usepackage{enumitem}
\usepackage{booktabs}
\usepackage{listings}
\usepackage{tcolorbox}
\usepackage{graphicx}
\usepackage{float}
\usepackage{amsmath,amssymb,bm}
\usepackage{xurl}

\hypersetup{pdftitle={Participation-Constrained Cell-Free ISAC: Project Wiki and Poster Defense Manual},pdfauthor={Xinpeng Yu}}
\setcounter{tocdepth}{2}
\graphicspath{{../../deliverables/paper_results_and_figures_v1_0/figures/}}
\lstset{basicstyle=\ttfamily\small,breaklines=true,frame=single,columns=fullflexible}
\newcommand{\mat}[1]{\mathbf{#1}}
\newcommand{\vect}[1]{\mathbf{#1}}
\newcommand{\quick}[1]{\begin{tcolorbox}[colback=blue!4,colframe=blue!55!black,title=Quick answer / 英文速答]\small #1\end{tcolorbox}}
\newcommand{\logic}[1]{\begin{tcolorbox}[colback=green!4,colframe=green!45!black,title=Chinese logic / 中文答题逻辑]\small #1\end{tcolorbox}}
\newcommand{\deep}[1]{\begin{tcolorbox}[colback=orange!5,colframe=orange!65!black,title=Follow-up detail / 追问扩展]\small #1\end{tcolorbox}}
\newcommand{\sourcebox}[1]{\begin{tcolorbox}[colback=gray!7,colframe=gray!65,title=Evidence and source]\raggedright\small #1\end{tcolorbox}}

\title{Participation-Constrained Cell-Free ISAC}
\subtitle{Project Wiki, Reproducibility Guide, and Poster Defense Manual}
\author{Xinpeng Yu}
\institute{UNSW Sydney --- School of Electrical Engineering and Telecommunications}
\date{30 July 2026}
\version{v1.0}
\extrainfo{English technical text with Chinese defense prompts.}
\cover{assets/solution_framework.png}

\begin{document}
\maketitle
\frontmatter
\tableofcontents

\chapter*{How to Use This Manual}
\addcontentsline{toc}{chapter}{How to Use This Manual}
Read Chapters 1--6 to understand and reproduce the project. Chapter 7 gives
spoken poster narratives. Chapter 8 is designed for rehearsal: state the blue
``Quick answer'' first, then use the Chinese logic and follow-up detail only
when the examiner asks for more depth. All reported numerical values are drawn
only from the final evidence package in
\texttt{deliverables/paper\_results\_and\_figures\_v1\_0}.

\mainmatter
\chapter{Project Overview and Research Claim}

\begin{introduction}
\item Explain the physical architecture before presenting the optimization.
\item Separate the mathematical relaxation from the certified physical solution.
\item Use the final evidence package as the sole source of numerical claims.
\end{introduction}

\section{Problem in one paragraph}
This project studies downlink resource allocation for a participation-constrained
Cell-Free integrated sensing and communication (ISAC) network. Distributed APs
jointly transmit communication streams to single-antenna UEs, while dedicated
target-specific sensing waveforms are emitted only by selected AP--target
pairs. The central processor minimizes total transmit power subject to robust
communication QoS, sensing-detection QoS, PCRB tracking accuracy, per-AP power,
and sensing-cluster cardinality constraints. Imperfect CSI and binary AP--target
participation make the design a nonconvex mixed-integer problem.

\section{Core contribution}
The contribution is not a global MINLP solver. It is a physically interpretable
and certifiable solution pipeline: covariance lifting and SDR expose an SDP
structure; robust SINR and PCRB constraints are represented by LMIs; double DC
penalties encourage rank-one and binary structure; and a separate recovery
stage fixes a binary topology, re-optimizes continuous covariances, and
validates every physical constraint.

\section{One-page oral summary}
\quick{We propose a Cell-Free ISAC design in which communication remains
globally cooperative, while dedicated sensing waveforms are authorized by a
target-centric AP cluster. The resulting robust mixed-integer design is handled
by SDR, the S-procedure, PCRB Schur LMIs, and double DC-SCA. A fixed-assignment
recovery and validation stage converts the relaxed result into a physically
feasible binary sensing solution.}
\logic{先讲架构的非对称性：通信全网协作，感知按目标选择集群。再讲难点：
鲁棒约束、PCRB 和二元关联耦合。最后强调贡献是\textbf{可认证的启发式恢复流程}，
不是声称原始混合整数问题的全局最优。}

\section{Claim boundaries}
\begin{itemize}
\item SDR is a continuous lower bound, not a deployable design.
\item A solver failure or time limit is not evidence of physical infeasibility.
\item Fig.~8 is a QoS-requirement sensitivity study, not a global Pareto frontier.
\item M12 results demonstrate scalability evidence and computational cost, not real-time operation.
\end{itemize}

\sourcebox{Primary sources: \texttt{PROBLEM\_FORMULATION.md},
\texttt{MASTER\_CONSTRAINTS.md}, and
\texttt{deliverables/paper\_results\_and\_figures\_v1\_0/README.md}.}

\chapter{Physical System Setting}

\section{Architecture and participation semantics}
The network has $M$ distributed APs, each with $N_t$ transmit antennas, $K$
single-antenna UEs, and $P$ sensing targets. A central processor coordinates
beamforming and AP--target sensing participation over fronthaul. Communication
covariances $\{\mat{W}_k\}$ are globally cooperative. Dedicated sensing
covariances $\{\mat{S}_p\}$ are target-specific.

\begin{figure}[H]
\centering
\includegraphics[width=0.82\linewidth]{fig1_system_architecture.png}
\caption{Asymmetric architecture: globally cooperative communication and
target-specific dedicated sensing participation.}
\end{figure}

The binary variable $b_{mp}\in\{0,1\}$ authorizes AP $m$ to emit nonzero
dedicated sensing power for target $p$. It is therefore an AP--target dedicated
sensing-participation indicator. It does \emph{not} gate AP $m$'s communication
transmission. With a nonzero sensing-power floor, an authorized pair has a
physical sensing contribution; without that floor, $b_{mp}=1$ should be read
only as authorization rather than proof of nonzero emission.

\section{Signal and power model}
Let $\vect{h}_k$ be the stacked AP-to-UE channel and let
$\mat{W}_k=\vect{w}_k\vect{w}_k^H\succeq0$ denote the covariance of UE $k$'s
communication beamformer. The total transmit covariance is
\[
  \mat{R}_X=\sum_{k=1}^{K}\mat{W}_k+\sum_{p=1}^{P}\mat{S}_p.
\]
Because a UE does not in general know the dedicated sensing waveform,
$\sum_p\mat{S}_p$ enters its interference term unless waveform cancellation is
explicitly assumed. Each AP obeys
$\operatorname{tr}(\mat{E}_m\mat{R}_X)\leq P_{\max}$, whereas dedicated sensing
participation is enforced through
\[
 \operatorname{tr}(\mat{E}_m\mat{S}_p)\leq P_{\max}b_{mp}.
\]

\section{Default operating point}
The principal study uses a $400\,\mathrm{m}\times400\,\mathrm{m}$ area,
$M=6$, $N_t=2$, $K=3$, $P=2$, $N_\theta=2$, and $P_{\max}=0.1$ W (20 dBm).
The robustness radius is $\epsilon_h=0.05$ unless scanned. The cluster size is
varied as $N_{\rm req}\in\{2,3,4,5,6\}$.

\section{MISO versus massive MIMO}
At the UE level, the current downlink model is multi-user MISO: each UE has one
receive antenna, while the AP network provides a distributed multi-antenna
transmit array. ``Cell-Free massive MIMO'' describes the distributed-array
architecture and its scalable extension; it does not imply that the present
UE receiver is a multi-antenna MIMO receiver. The implementation deliberately
uses the MISO-equivalent observation model rather than a Kronecker MIMO
observation model requiring a different receiver architecture.

\sourcebox{Implementation sources: \texttt{sim/matlab/generate\_scenario.m},
\texttt{default\_params.m}, and \texttt{build\_E\_m.m}.}

\chapter{Mathematical Model}

\section{Robust communication QoS}
For an uncertain channel $\vect{h}_k=\hat{\vect{h}}_k+\Delta\vect{h}_k$ with
$\|\Delta\vect{h}_k\|_2\leq\epsilon_h$, the worst-case communication SINR
requirement is imposed over the entire uncertainty ball. In covariance form,
the desired communication covariance must dominate multiuser interference,
dedicated sensing interference, and noise for every admissible channel error.
This is a semi-infinite constraint before the S-procedure is applied.

\section{Sensing SINR and PCRB}
The dedicated covariance $\mat{S}_p$ is used to guarantee target detection
through a sensing-SINR threshold. Tracking accuracy is controlled by the data
FIM $\mat{J}_p$ and the trace-PCRB condition
\[
 \operatorname{tr}(\mat{J}_p^{-1})\leq\Gamma_{\rm track,p}.
\]
For the adopted complex Gaussian mean model, the data FIM contains the factor
$2/\sigma_s^2$. The physical derivative matrix $\mat{D}_p$ retains the target
geometry. A small numerical regularization may be used only to prevent singular
linear algebra; it must not be misrepresented as new sensing information.

\section{Original problem P1}
The original formulation jointly minimizes total transmit power over
$\{\vect{w}_k\}$, $\{\mat{S}_p\}$, and $\{b_{mp}\}$ subject to robust
communication QoS, sensing SINR, PCRB, per-AP power, sensing authorization, and
cardinality constraints. The coupled sources of nonconvexity are: beamformer
quadratic products, rank-one covariance relations, worst-case CSI constraints,
PCRB inverse, SINR ratios, and binary variables.

\section{Interpretation of PCRB geometry}
Sensing SINR is largely an energy measure, whereas PCRB depends on the
eigenstructure of $\mat{J}_p$. APs with similar viewing geometry can contribute
energy without providing a second informative direction. This explains why a
distance-only association can be poor for tracking, and why FIM-aware topology
construction is a strong baseline.

\sourcebox{Full derivation reference: \texttt{PROBLEM\_FORMULATION.md},
Sections 2--6; numerical implementation: \texttt{solve\_p3\_sca\_t.m}.}

\chapter{From P1 to the Dual DC-SCA Subproblem}

\section{P2: covariance lifting and SDR}
Introduce $\mat{W}_k=\vect{w}_k\vect{w}_k^H$. The equality is equivalent to
$\mat{W}_k\succeq0$ plus $\operatorname{rank}(\mat{W}_k)=1$. Dropping the rank
constraint yields an SDP relaxation. This is a relaxation, not an assertion
that rank one is automatically guaranteed in the present coupled robust ISAC
problem.

\section{Robust and PCRB LMIs}
The S-procedure converts the norm-bounded worst-case communication constraint
to an LMI with a nonnegative multiplier. The trace-PCRB constraint is expressed
through a Schur block LMI using an auxiliary matrix $\mat{M}_p$ and
$\operatorname{tr}(\mat{M}_p)\leq\Gamma_{\rm track,p}$. The full state dimension
$N_\theta$ is retained in this block; reducing it to an unrelated scalar block
would change the model.

\section{P3: double DC-SCA}
The rank penalty is based on the gap between trace and dominant eigenvalue; the
binary penalty is based on $b_{mp}(1-b_{mp})$. At each iteration, concave terms
are linearized around the previous iterate, producing a convex SDP subproblem.
The approach is a stationary-point heuristic with penalty continuation, not a
global certificate for the original mixed-integer problem.

\begin{figure}[H]
\centering
\includegraphics[width=0.82\linewidth]{fig5_statistical_double_dc_convergence.png}
\caption{Continuous double-DC stabilization over 25 common realizations.
The figure supports support stabilization before discrete recovery; it is not
itself a physical-feasibility result.}
\end{figure}

\section{Why the recovery step is separate}
The stable Top-$N_{\rm req}$ support of a continuous relaxation suggests a
binary topology but does not prove it physically feasible. Therefore, the
algorithm projects to binary $b$, solves a fixed-assignment continuous problem,
extracts rank-one beamformers when needed, and runs an explicit validator. This
separation is scientifically important: it prevents fractional relaxation
quality from being reported as a deployable result.

\sourcebox{Derivation source: \texttt{CONVEXIFICATION\_CHAIN.md}; caveats:
\texttt{SDR\_TIGHTNESS.md} and \texttt{COMPLEXITY\_BOUND.md}.}

\chapter{Algorithm, Code, and Reproducibility}

\section{End-to-end pipeline}
\begin{enumerate}[label=\textbf{Step \arabic*:},leftmargin=2.4cm]
\item Generate AP, UE, and target geometry; calibrate observable PCRB thresholds.
\item Solve the continuous SDR/DC-SCA subproblems and record rank, binary, and support diagnostics.
\item Construct binary candidates: stable Top-$N_{\rm req}$, FIM-greedy, and repair candidates when required.
\item For each candidate, solve the fixed-$b$ robust covariance problem.
\item Extract implementable beamformers when needed and validate cardinality, binary status, power, robust SINR, sensing SINR, PCRB, PSD, and LMI residuals.
\end{enumerate}

\section{Primary MATLAB map}
\begin{longtable}{p{0.30\linewidth}p{0.61\linewidth}}
\toprule
File & Responsibility\\ \midrule
\texttt{baseline\_alg2.m} & Orchestrates SDR/DC-SCA, candidate construction, recovery, and final status.\\
\texttt{solve\_p3\_sca\_t.m} & Builds and solves one convex dual-DC SCA subproblem.\\
\texttt{solve\_p3\_with\_fixed\_b.m} & Re-optimizes continuous covariances after a binary topology is fixed.\\
\texttt{validate\_solution.m} & Independently checks the returned physical solution.\\
\texttt{generate\_scenario.m} & Geometry, channels, derivatives, and automatic PCRB calibration.\\
\texttt{evaluate\_isac\_metrics.m} & Recomputes power, rates, PCRB, and sensing metrics.\\
\texttt{run\_*.m} and \texttt{plot\_*.m} & Reproducible Monte Carlo campaigns and figures.\\
\bottomrule
\end{longtable}

\section{What counts as a physical solution}
An accepted record must have a binary association, exact target cardinality,
authorized dedicated sensing power, per-AP power within tolerance, PSD
covariances, nonnegative S-procedure multipliers, robust LMI satisfaction, and
all QoS/PCRB checks passed. A fractional relaxation, an unverified
principal-eigenvector extraction, or a solver status alone is insufficient.

\section{Reproducibility protocol}
Use deterministic seed lists, common scenarios for method comparisons, and
unconditional denominators for feasibility. Condition power, sum rate, PCRB,
and sensing SINR on physically feasible solutions unless a common-feasible set
is stated. Record MATLAB, CVX, MOSEK, machine, solver tolerance, seed list, and
Git commit with every campaign.

\sourcebox{Authoritative frozen scripts and data snapshots are in
\texttt{deliverables/paper\_results\_and\_figures\_v1\_0/scripts} and
\texttt{.../results}. The active development copies are under \texttt{sim/matlab}.}

\chapter{Results, Evidence, and Safe Interpretation}

\section{Evidence rules}
The final evidence package is the sole numerical authority. Feasibility is an
unconditional rate over all trials at the stated point. Other performance
statistics are conditional on physical feasibility unless explicitly described
as paired or common-feasible. These rules avoid silently discarding hard cases.

\section{Core results}
\begin{figure}[H]
\centering
\includegraphics[width=0.93\linewidth]{fig3_cluster_size_tradeoff.png}
\caption{Cluster-size trade-off: feasibility, conditional total transmit power,
and runtime. The lowest observed mean power in the M6 study occurs at an
intermediate sensing-cluster size.}
\end{figure}

Across 30 common M6 scenarios, the proposed method is physically feasible at
each tested $N_{\rm req}\in\{2,3,4,5,6\}$. Mean power is 37.65, 30.83, 29.50,
29.74, and 30.94 mW, respectively. The nonmonotone relation is meaningful:
small clusters lack geometric FIM diversity, while large clusters incur
dedicated-waveform resource and interference costs.

\begin{figure}[H]
\centering
\includegraphics[width=0.93\linewidth]{fig10_method_comparison_vs_nreq.png}
\caption{Association/recovery comparison. FIM-aware selection is a strong
topology baseline; full recovery performs the final coupled robust covariance
optimization and validation.}
\end{figure}

At $N_{\rm req}=3$, mean power is 30.83 mW for the proposed method, 30.93 mW
for FIM-greedy, and 39.29 mW for nearest-AP. The close FIM-greedy result shows
that geometry strongly selects the sensing topology; it does not eliminate the
need for robust beamforming, fixed-topology re-optimization, and certification.

\section{Robustness and QoS sensitivity}
\begin{figure}[H]
\centering
\includegraphics[width=0.86\linewidth]{fig9_csi_robustness.png}
\caption{Sampled CSI robustness. The Monte Carlo result complements, but does
not replace, the worst-case S-procedure formulation.}
\end{figure}

The robust design has zero sampled system outage over the reported uncertainty
radii; the nominal-CSI design has mean outage around 89--92\%. Fig.~8 should be
called a \emph{QoS-requirement sensitivity study}: stricter communication and
tracking targets require more power over the tested grid. It is neither a
global communication--sensing Pareto frontier nor the boundary of the full
feasible region.

\section{Scalability and limitations}
The M12, $K=6$, $P=3$ fixed-cardinality study reports proposed recovery feasible
in 8/8 cases at $N_{\rm req}=3$. Its mean end-to-end runtime is 496.3 seconds
per scenario. This is evidence of larger-system applicability and a clear
motivation for future acceleration, not a real-time claim.

\sourcebox{Detailed bilingual figure explanations and all precise values are in
\texttt{RESULTS\_AND\_FIGURES\_BILINGUAL.md}.}

\chapter{Poster Presentation Scripts}

\section{30-second version}
\quick{This work studies robust resource allocation for Cell-Free ISAC. We keep
communication globally cooperative, but activate dedicated sensing waveforms
only for a target-centric AP cluster. The resulting mixed-integer robust design
is solved through SDR, robust LMIs, double DC-SCA, and a certified recovery
stage. The key result is that FIM-aware sensing clusters improve feasibility and
energy efficiency compared with distance-only association.}
\logic{30 秒只讲问题、架构非对称性、方法链和一个结果。不要讲所有约束，不要先讲
推导细节。}

\section{2-minute version}
\quick{Cell-Free ISAC must serve UEs and track targets with the same distributed
AP network. Our model separates these roles: every AP may cooperate for
communication, while the binary variable $b_{mp}$ authorizes only dedicated
sensing transmission for target $p$. We minimize power under worst-case
communication SINR, sensing SINR, PCRB tracking, and per-AP constraints. After
covariance lifting, we use the S-procedure and Schur LMIs, then a double DC-SCA
algorithm to promote rank-one covariances and binary associations. Because a
relaxed result is not physical by itself, we project a stable support, solve a
fixed-assignment problem, and validate it. Experiments show an intermediate
cluster size is most energy efficient, robust design removes sampled outage,
and FIM-aware recovery outperforms nearest-AP association in difficult geometry.}
\logic{两分钟版本的顺序是：\textbf{architecture → problem → solution → certification → evidence}。
看到老师点头后停在这里，把时间留给提问。}

\section{5-minute version outline}
\begin{enumerate}
\item Motivation: shared AP infrastructure creates communication--sensing resource coupling.
\item Architecture: global $\mat{W}_k$ versus target-specific $\mat{S}_p$ and $b_{mp}$.
\item Constraints: robust SINR, sensing SINR, PCRB, power, and cardinality.
\item Method: P1 to P2 lifting/SDR, P3 double DC-SCA, then certified recovery.
\item Results: binary stabilization, $N_{\rm req}$ trade-off, baseline comparison, robustness, and M12 cost.
\item Closing: geometric sensing selection plus robust continuous optimization; future distributed acceleration.
\end{enumerate}

\section{Transitions worth memorizing}
\begin{itemize}
\item ``The important modeling choice is that $b_{mp}$ limits dedicated sensing, not global communication.''
\item ``The relaxation identifies a promising support; physical feasibility is established only after fixed-topology recovery.''
\item ``This figure is a QoS sensitivity study, not a Pareto-boundary claim.''
\item ``The large-scale runtime is a limitation we report transparently, not a real-time performance claim.''
\end{itemize}

\chapter{Poster Defense Q\&A Bank}

\newcommand{\defenseqa}[5]{\subsection{Q#1. #2}\quick{#3}\logic{#4}\deep{#5}}

\section{Research motivation and contribution}
\defenseqa{1}{What problem does this project solve?}{It minimizes transmit power in a Cell-Free ISAC network while guaranteeing robust communication, sensing detection, tracking accuracy, AP power limits, and sparse sensing participation.}{一句话是：同一套分布式 AP 同时服务用户和跟踪目标，资源耦合导致优化困难。}{The novelty is the participation-constrained dedicated sensing architecture plus a physically validated recovery pipeline. See Chapters 1--4.}

\defenseqa{2}{Why use Cell-Free ISAC?}{Cell-Free cooperation removes fixed cell boundaries and lets distributed APs jointly shape communication and sensing resources according to geometry.}{强调分布式 AP 的空间分集，尤其适合目标定位而不是只增加通信信号强度。}{The benefit is conditional on fronthaul coordination and synchronization; those costs are listed as future work.}

\defenseqa{3}{What is the central modeling idea?}{Communication is globally cooperative, whereas dedicated sensing is target-centric and sparse. This asymmetry preserves communication degrees of freedom while limiting sensing processing and waveform resources.}{核心关键词：global communication, local sensing。}{Do not say that unselected APs are switched off; they may still transmit communication data.}

\defenseqa{4}{What is new beyond conventional AP association?}{The association controls dedicated sensing-waveform authorization and is optimized jointly with robust transmit covariances and PCRB constraints.}{不是传统只看距离或信道增益的用户关联；PCRB 需要空间几何多样性。}{FIM-greedy is a strong geometry baseline, but it does not replace coupled robust covariance optimization.}

\defenseqa{5}{Why minimize power?}{Power is a common system cost that directly exposes the price of robust communication, sensing, and tracking requirements.}{功耗是统一资源指标，便于比较不同拓扑和 QoS 压力。}{The objective is total transmit power; fronthaul and computation energy are not yet included.}

\defenseqa{6}{Is this a global MINLP optimum?}{No. The method is a structured relaxation, DC-SCA, and recovery heuristic that returns validated feasible solutions, not a global optimality certificate.}{必须主动承认：非凸混合整数问题没有全局最优保证。}{The SDR value is used only as a lower bound for interpreting a relaxation gap.}

\defenseqa{7}{What is the main empirical finding?}{FIM-aware sensing topology and certified recovery improve feasibility and energy efficiency relative to nearest-AP and random association, particularly in difficult geometry.}{回答时给一条数字：$N_{\rm req}=3$ 下 proposed 30.83 mW，nearest-AP 39.29 mW。}{The result is backed by common-seed Monte Carlo comparisons in the final evidence package.}

\defenseqa{8}{What should the audience remember?}{Tracking quality depends on sensing geometry, not only received energy; therefore a physically validated, FIM-aware sensing cluster is essential in Cell-Free ISAC.}{记忆点：PCRB 看 FIM 特征值和几何，不只是功率。}{Use this as the closing sentence after discussing the experiments.}

\section{Physical model and architecture}
\defenseqa{9}{What exactly does $b_{mp}$ mean?}{$b_{mp}=1$ authorizes AP $m$ to transmit nonzero dedicated sensing power for target $p$; it is an AP--target sensing-participation indicator.}{它是感知任务参与/授权变量，不是通信 AP 开关。}{With a sensing-power floor, it represents actual nonzero sensing participation; otherwise it is only an upper-bound authorization.}

\defenseqa{10}{Does $b_{mp}$ gate communication transmission?}{No. Communication covariance is globally cooperative and is not gated by $b_{mp}$.}{这是海报最容易被问到的点，回答必须坚定。}{The dedicated sensing covariance obeys the Big-M authorization constraint, while total AP power includes both communication and sensing.}

\defenseqa{11}{Why is dedicated sensing interference included at UEs?}{A UE cannot generally cancel an unshared dedicated sensing waveform, so its received energy is modeled as interference in the communication SINR denominator.}{如果波形未知，UE 无法 SIC，因此必须算干扰。}{A cancellation model would be a different physical architecture and would require explicit receiver-side assumptions.}

\defenseqa{12}{Why not make every AP sense every target?}{Using every AP can improve diversity but consumes dedicated waveform, power, processing, synchronization, and fronthaul resources.}{全参与不是免费；$N_{\rm req}$ 控制感知规模。}{The measured power relation is nonmonotone because diversity gains eventually saturate while sensing costs persist.}

\defenseqa{13}{Why use dedicated sensing covariance $\mat{S}_p$?}{It gives each target an explicitly controlled sensing waveform and makes AP--target participation physically meaningful.}{这样 $b_{mp}$ 真正控制专用感知资源，而不是空洞标签。}{Communication beams are intentionally not credited toward PCRB in this partially shared signaling model.}

\defenseqa{14}{Is the system MIMO or MISO?}{The implemented downlink is multi-user MISO: UEs are single-antenna and the distributed AP network is the multi-antenna transmitter.}{不要把 Cell-Free massive MIMO 的架构名称误说成 UE 端 MIMO 接收机。}{The M12 experiment scales the distributed transmit dimension, not the UE receive dimension.}

\defenseqa{15}{Why not use a Kronecker MIMO observation model?}{That model would imply a more complex receive architecture and would no longer match the implemented MISO-equivalent signal model.}{理论、代码和接收机假设必须闭环；不能为了公式复杂而改变物理模型。}{The manuscript should state the MISO-equivalent observation assumption explicitly.}

\defenseqa{16}{How is AP power enforced?}{Each AP has a hard per-AP constraint on the corresponding block of the total covariance $\mat{R}_X$.}{功率单位统一使用瓦特，$0.1$ W 等于 20 dBm。}{This prevents a relaxed solution from satisfying QoS by concentrating unbounded power at one AP.}

\defenseqa{17}{What is the role of $N_{\rm req}$?}{It fixes the number of APs participating in the dedicated sensing cluster for each target and is a key system design variable.}{它控制感知的空间自由度和资源成本。}{It is swept from 2 to 6 in the main M6 study and from 2 to 5 in the M12 study.}

\section{Mathematical formulation}
\defenseqa{18}{Why is the communication SINR robust?}{The design guarantees the required SINR for every channel error in a bounded uncertainty ball, rather than only at the estimated channel.}{讲清楚：CSI 有误差，所以不是只在 $\hat h$ 上满足。}{The S-procedure is exact for the single quadratic implication under its standard regularity condition.}

\defenseqa{19}{What does the PCRB constraint measure?}{It bounds the trace of the inverse target FIM, which summarizes the achievable estimation-error covariance for the target-state parameters.}{PCRB 是跟踪精度，不等同于检测 SINR。}{A small sensing SINR alone does not ensure a well-conditioned FIM.}

\defenseqa{20}{Why can sensing SINR and PCRB disagree?}{Sensing SINR measures useful energy, while PCRB also requires independent geometric information directions in the FIM.}{同方向的高能量观测可能仍然定位差。}{This is why FIM geometry is used in topology construction.}

\defenseqa{21}{Why is there a factor of two in the FIM?}{For the adopted proper complex Gaussian mean model, the real FIM contribution includes $2/\sigma_s^2$.}{这是复高斯均值模型的标准系数；漏掉会让 PCRB 标度错误。}{The implementation and metric evaluator use the same convention.}

\defenseqa{22}{How is the PCRB inverse handled?}{A Schur-complement LMI introduces an auxiliary matrix so that the trace-PCRB bound is represented without directly inverting a variable matrix.}{回答结构：inverse 很难，Schur LMI 把它变成凸约束。}{The block retains the full $N_\theta$ state dimension; replacing it by a scalar would be a different approximation.}

\defenseqa{23}{What makes P1 nonconvex?}{Rank-one beamforming, binary participation, SINR ratios, worst-case uncertainty, and the inverse-FIM PCRB constraint are jointly nonconvex.}{不要只说 binary；rank、robust、PCRB 都是难点。}{The convexification chain documents each source and its treatment.}

\defenseqa{24}{What is covariance lifting?}{It replaces beamforming-vector products by PSD covariance matrices, making received powers affine in the covariance variables.}{$W_k=w_kw_k^H$ 后，二次项变 trace 线性项。}{The price is an explicit rank-one constraint, which is subsequently relaxed and penalized.}

\defenseqa{25}{What is SDR?}{Semidefinite relaxation drops the rank-one constraint and solves a convex PSD covariance problem.}{SDR 给下界和初始化，但不是最终可部署方案。}{Rank tightness is empirical here; it is not claimed as a general theorem.}

\defenseqa{26}{What is the S-procedure doing?}{It converts a norm-bounded, infinitely many-channel robust SINR condition into one finite LMI with a nonnegative multiplier.}{无穷多个 CSI 扰动约束变为一个矩阵不等式。}{The multiplier must also be checked nonnegative by the validator.}

\defenseqa{27}{Why use trace-PCRB rather than one scalar FIM gain?}{The trace of the inverse FIM penalizes weak state directions and therefore reflects multi-dimensional tracking quality.}{不能只看 FIM 总能量；最小特征值太小会主导误差。}{D-optimal or E-optimal scores are used only as topology heuristics, not as replacements for the PCRB constraint.}

\section{Optimization and recovery}
\defenseqa{28}{What is the purpose of the rank DC penalty?}{It encourages a PSD covariance to approach rank one by penalizing the gap between its trace and dominant eigenvalue.}{秩一意味着可从协方差恢复单流波束。}{In the current MU-MISO tests, SDR is already nearly rank one, so the observed dominant benefit is binary recovery.}

\defenseqa{29}{What is the binary DC penalty?}{It penalizes the concave binary residual $b_{mp}(1-b_{mp})$ while maintaining a convex linearized subproblem at each iteration.}{$b(1-b)$ 在 0 和 1 为零，中间最大。}{The reported median binary residual drops from roughly $4.3\times10^{-1}$ to about $6\times10^{-5}$.}

\defenseqa{30}{Why call it double DC-SCA?}{Two nonconvex structural penalties, rank and binary, are successively linearized within an SCA framework.}{double 指两类 DC 结构，不是两个独立算法。}{The robust and PCRB LMIs are already convex after reformulation.}

\defenseqa{31}{Does DC-SCA converge?}{The implemented procedure shows rapid empirical stabilization of support and residuals; it does not prove global convergence of P1.}{用“stabilization”而不是“global convergence”。}{The 25-seed trace shows support stability from the first DC iteration and binary residual near $10^{-4}$ by the second.}

\defenseqa{32}{Why stop after support stabilization?}{Once the Top-$N_{\rm req}$ support stops changing, further continuous iterations often do not alter the discrete candidate, so recovery can begin.}{节省 SDP 时间，但最后必须做 fixed-$b$ 验证。}{The stopping signal is not a replacement for physical feasibility.}

\defenseqa{33}{Why is direct rounding insufficient?}{A small fractional residual can still represent a topology that becomes infeasible after exact binary projection.}{连续解可借用很小的分数自由度平衡严格约束。}{The recovery ablation shows direct DC Top-N rounding has much lower feasibility than full recovery.}

\defenseqa{34}{What happens in fixed-$b$ re-optimization?}{The association is held binary and the continuous communication and sensing covariances are re-optimized under all original physical constraints.}{固定拓扑后重新找最好的波束与功率，不能直接复用松弛协方差。}{This is the key bridge from a promising relaxed support to a certified physical solution.}

\defenseqa{35}{Why include FIM-greedy candidates?}{They cheaply construct geometry-diverse sensing clusters that target informative FIM directions and repair poor DC rounding choices.}{FIM-greedy 是纯代数、速度快、物理直觉强。}{It is a candidate generator and strong baseline, not an oracle or global combinatorial solver.}

\defenseqa{36}{What is validated after recovery?}{Binary/cardinality constraints, authorized sensing power, per-AP power, PSD, robust LMIs, PCRB, communication SINR, and sensing SINR are checked independently.}{强调“独立验证器”，避免只相信 CVX status。}{A solver status of Infeasible or Failed is recorded separately from a validated physical infeasibility conclusion.}

\section{Code and reproducibility}
\defenseqa{37}{Where is the main algorithm implemented?}{The development implementation is centered on \texttt{sim/matlab/baseline\_alg2.m}, with separate SCA, fixed-$b$, scenario, metric, and validation functions.}{主入口不是一个文件完成所有事，而是模块化 pipeline。}{The frozen scripts used for final evidence are copied under the final deliverables package.}

\defenseqa{38}{Why keep a separate validator?}{It prevents the optimization and reporting code from sharing the same modeling omission unnoticed.}{独立验证是防御性科研实践。}{For example, it checks the sensing-interference term, multiplier signs, binary association, and cardinality explicitly.}

\defenseqa{39}{How are random experiments made comparable?}{Competing methods use identical common seeds and geometry within each trial.}{同一 seed 下比较才可归因于方法差异，不是场景差异。}{Conditional power comparisons are reported only on feasible cases or paired common-feasible sets.}

\defenseqa{40}{What solver stack is used?}{The recorded final environment uses MATLAB R2026a, CVX 2.2, and MOSEK 11.2.2.}{版本要和证据包 README 一致，勿混用旧元数据。}{The build manual records the command and warns that solver versions must be frozen before submission.}

\defenseqa{41}{How are units handled?}{Power and noise variances are consistently represented in linear watts; for example, 20 dBm is 0.1 W.}{单位不一致会造成数千倍 SNR 标度错误。}{All input parameters and evaluation metrics should be checked in the same linear unit system.}

\defenseqa{42}{How is automatic PCRB calibration protected?}{The reference FIM is checked for observability and conditioning before using its inverse to set a threshold.}{若参考 FIM 奇异，不能把 Inf 或 NaN 当作正常阈值。}{A calibration rejection is distinct from a solver-declared infeasibility.}

\defenseqa{43}{What is the reproducibility minimum?}{A result needs the exact seed list, configuration, code commit, MATLAB/CVX/MOSEK versions, machine information, and saved MAT/CSV records.}{只有图不够；必须能从数据和脚本复画。}{The final evidence package stores figures, raw summaries, scripts, and reporting conventions.}

\section{Experiments and interpretation}
\defenseqa{44}{Why does an intermediate $N_{\rm req}$ minimize power?}{More APs initially add geometric FIM diversity, but beyond a point dedicated sensing waveform and interference costs outweigh the diversity gain.}{这是“少了不够信息，多了资源浪费”的平衡。}{In M6, the observed minimum mean power occurs at $N_{\rm req}=4$.}

\defenseqa{45}{Why is nearest-AP a weak sensing baseline?}{Distance alone ignores derivative directions and angular diversity, which determine the conditioning of the PCRB FIM.}{最近不代表从不同角度看目标。}{Nearest-AP remains a useful theoretical geometry reference, not an online policy requiring unknown exact target positions.}

\defenseqa{46}{Why is FIM-greedy close to the proposed method?}{The sensing topology is strongly governed by geometry, so a good FIM score selects many of the same APs as the optimized recovery.}{这不是坏事，说明优化发现了有物理意义的几何规律。}{The proposed method still performs the coupled robust covariance optimization and validation that FIM-greedy alone does not provide.}

\defenseqa{47}{How should Fig. 5 be described?}{It is a 25-realization statistical stabilization trace for the continuous double-DC phase before discrete recovery.}{不要说“每轮都产生物理解”。}{It shows support stabilization and small residuals; fixed-$b$ recovery remains mandatory.}

\defenseqa{48}{How should Fig. 8 be described?}{It is a QoS-requirement sensitivity study: the required mean transmit power is mapped over communication SINR and PCRB allowance targets.}{不要叫 Pareto frontier。}{All tested grid points are feasible, but the figure does not identify the entire feasible-region boundary.}

\defenseqa{49}{What does zero sampled outage prove?}{It shows strong empirical robustness over the sampled channel perturbations, consistent with the worst-case design.}{先说 sampled，再说 worst-case certificate；不要把样本当完整证明。}{The S-procedure constraint, not Monte Carlo sampling alone, is the formal robustness mechanism.}

\defenseqa{50}{Why is nominal outage near 90 percent?}{A design optimized only for the estimated channel can violate a tight SINR target under even moderate perturbations.}{说明鲁棒约束有实际价值，而非纯数学复杂化。}{The robust power premium at the reported high uncertainty point remains modest relative to the outage reduction.}

\defenseqa{51}{Are PCRB and SINR constraints actually active?}{The reported QoS-tightness checks show PCRB ratios essentially at their boundaries and nonnegative communication and sensing margins.}{防止别人怀疑低功耗是因为偷偷放松约束。}{Metrics are recomputed from returned covariances, not taken only from solver internals.}

\defenseqa{52}{What does the M12 study add?}{It tests a larger network with $M=12$, $K=6$, and $P=3$, showing the method can still find validated solutions while exposing substantial runtime.}{大规模验证既展示可扩展性，也诚实展示代价。}{At fixed $N_{\rm req}=3$, the proposed method is feasible in 8/8 reported cases.}

\defenseqa{53}{Why are some statistics conditional on feasibility?}{Power or rate is undefined for a method that has no validated physical solution, so those values are conditioned on feasible cases while feasibility itself remains unconditional.}{先报成功率，再比较成功样本性能。}{For fair method gaps, use paired common-feasible sets when explicitly stated.}

\section{Limitations and future work}
\defenseqa{54}{What is the largest current limitation?}{The robust SDP and recovery pipeline is computationally expensive, especially at larger transmit dimension.}{不要回避 496 秒；把它作为下一步动机。}{Distributed, first-order, warm-started, or learning-assisted methods are natural acceleration directions.}

\defenseqa{55}{Does the model include fronthaul capacity?}{No. The current model interprets sparse sensing participation as a processing/fronthaul-motivated abstraction, but it does not yet optimize an explicit fronthaul capacity constraint.}{这是合理但必须承认的模型边界。}{A future extension can add fronthaul rate, synchronization, and processing-energy terms.}

\defenseqa{56}{Does the model include clutter and hardware impairments?}{Not explicitly. It uses the stated channel, noise, CSI-ball, and FIM assumptions as an analytical baseline.}{承认理想化假设，但说明鲁棒 CSI 是迈向实践的一步。}{Future work includes clutter, calibration error, waveform nonidealities, and hardware-in-the-loop testing.}

\defenseqa{57}{What if a target moves?}{The present optimization is a static snapshot; a dynamic system would repeatedly update geometry and perform cluster handover or predictive warm starts.}{不要说当前已解决动态问题。}{The stable-support observation motivates incremental updates, but dynamic handover is future work.}

\defenseqa{58}{Can machine learning replace the optimizer?}{Learning can accelerate candidate initialization or predict promising clusters, but the final design still needs optimization-based physical constraint certification.}{AI 适合 warm start，不应替代可行性认证。}{A hybrid approach can retain fixed-$b$ validation as a safety layer.}

\defenseqa{59}{Why not claim real-time applicability?}{The current centralized SDP runtime is too high for a real-time claim; the contribution is a rigorous reference design and evidence base for faster future methods.}{诚实比夸张更重要。}{The M12 runtime is explicitly reported as computational cost.}

\defenseqa{60}{What is the next concrete research step?}{Develop a distributed or first-order solver with topology warm starts, then validate it under moving targets, fronthaul limits, and realistic radar/channel measurements.}{未来工作要具体：算法加速、动态、物理非理想、实验验证。}{This preserves the current model's strongest feature: certified coupling of sensing topology and robust covariances.}


\appendix
\chapter{Poster Consistency Ledger}

This ledger is deliberately concise. It records statements that must remain
consistent across the poster, manuscript, code, and oral answers.

\begin{longtable}{p{0.30\linewidth}p{0.61\linewidth}}
\toprule
Topic & Approved statement\\ \midrule
Receiver model & The current implementation is a Cell-Free multi-user MISO downlink with single-antenna UEs and distributed AP transmit arrays.\\
Participation $b_{mp}$ & It authorizes dedicated sensing transmission for target $p$; it does not gate globally cooperative communication.\\
P2 & Covariance lifting introduces rank-one constraints, and SDR relaxes them. P2 does not preserve an explicit rank-one convex constraint.\\
P3 & P3 denotes the convex DC-SCA subproblem solved at each iteration, not the original MINLP.\\
Step 4 & Certified physical recovery is an algorithmic step: project support, fix $b$, re-optimize, extract if needed, and validate. It is not a same-level theoretical P4.\\
SDR & SDR is a lower bound / relaxation reference, never a deployable physical baseline.\\
Fig. 5 & A continuous-phase stabilization diagnostic before discrete recovery.\\
Fig. 8 & A QoS-requirement sensitivity study, not a global Pareto frontier.\\
M12 runtime & A computational cost measurement, not evidence of real-time feasibility.\\
Global optimality & Not claimed. The reported method is a validated heuristic recovery pipeline for a nonconvex MINLP.\\
\bottomrule
\end{longtable}

\section{Poster-specific wording repairs}
Use ``Step 1--4'' consistently in the solution framework. Keep P1, P2, and P3
in the first three card titles to match the derivation. Call the final card
``Certified Physical Recovery,'' not P4. Use ``QoS-requirement sensitivity''
for the Fig.~8 description. Refer to the method as ``Cell-Free ISAC'' and,
when challenged about receiver dimensionality, state ``MU-MISO downlink
equivalent model.''

\chapter{Build and Reproduction Guide}

\section{Build the manual}
The manual is designed for XeLaTeX because it includes Chinese defense prompts.
From \texttt{docs/project\_wiki}, run:
\begin{lstlisting}[language=bash]
latexmk -xelatex -interaction=nonstopmode main.tex
\end{lstlisting}
The vendored ElegantBook class is resolved through the document-class path, so
no system-wide \texttt{elegantbook} installation is required.

\section{Expected environment}
\begin{itemize}
\item TeX Live 2025 or later with XeLaTeX, CTeX, biblatex, and tcolorbox.
\item MATLAB R2026a, CVX 2.2, and MOSEK 11.2.2 for the recorded experiments.
\item The local final evidence package at
\texttt{deliverables/paper\_results\_and\_figures\_v1\_0}.
\end{itemize}

\section{Troubleshooting}
\begin{longtable}{p{0.34\linewidth}p{0.57\linewidth}}
\toprule
Symptom & Action\\ \midrule
Chinese glyphs are missing & Use XeLaTeX, not pdfLaTeX; install a complete CTeX distribution.\\
Figures are missing & Confirm the final evidence package exists at the repository-relative path; do not substitute exploratory figures.\\
Undefined references & Run \texttt{latexmk} again, or delete only generated files with \texttt{latexmk -C}.\\
ElegantBook class unavailable & Do not install a random package version; use the bundled vendor snapshot.\\
MATLAB values differ & Check seed list, solver version, $P_{\max}$ unit, $\Gamma$ calibration, and whether metrics are conditioned on physical feasibility.\\
\bottomrule
\end{longtable}

\section{Final validation checklist}
Before distribution, compile twice, verify that the PDF has no unresolved
references, render representative pages, and inspect the title page, long
tables, formulas, figure captions, and the Q\&A chapter. Check that every
numerical statement can be found in the final evidence package and that the
Git commit records only files under \texttt{docs/project\_wiki}.

\backmatter
\end{document}
